% Anexo: los cálculos de la solución, reunidos (con los datos del cuadernillo)
\clearpage
\section*{Anexo: los cálculos de la solución}

\begin{itemize}
\item Todos los cálculos que usa la solución salen de los cuadros y gráficos del cuadernillo con las fórmulas de la guía (A.2 y A.3). Aquí se reúnen en una tabla, con el dato, la operación y la pregunta en la que se usan; el detalle está en los cuadros complementarios de cada pregunta.
\end{itemize}

\begin{table}[H]
\centering\scriptsize
\renewcommand{\arraystretch}{1.12}
\caption{Los cálculos de la solución}\label{tab:anexo-calculos}
\begin{tabular}{>{\raggedright\arraybackslash}p{3.6cm}>{\raggedright\arraybackslash}p{6.9cm}>{\centering\arraybackslash}p{2.3cm}>{\centering\arraybackslash}p{1.6cm}}
\toprule
\rowcolor{tceeazulpalido}
\textbf{Cálculo} & \textbf{Operación con los datos del cuadernillo} & \textbf{Resultado} & \textbf{Pregunta} \\
\midrule
\rowcolor{tceesalvia}\multicolumn{4}{l}{\emph{Ciclo y demanda}} \\
Tasa anualizada, 2T2026 & $(1+0{,}0068)^4-1$ (G1) & 2,8~\% & 1 \\
PIB medio de 2026 & niveles encadenados de G1; 3T y 4T al 0,4~\% & $\approx2{,}5~\%$ & 1 \\
Potencial implícito, 2025 & $\gamma_Y-\Delta\mathrm{OG}=2{,}6-(0{,}6-0{,}5)$ (C1) & 2,5~\% & 1 \\
PIB y PIB por habitante, 2019-2025 & $118{,}7/107{,}9$ y $108{,}7/103{,}6$ (G2) & +10,0 y +5,0~\% & 1 \\
Descomposición del PIB, 2019-2025 & $0{,}79-0{,}01+0{,}95-0{,}46+0{,}33$ (C2 bis) & 1,59 & 1 \\
FBCF acumulada por activos & $\prod(1+\gamma_t)-1$ (C1) & intangibles +47,4~\% (2019-2025) & 1 \\
Deuda de los hogares en euros & $0{,}565\times1.254$ y $0{,}429\times1.690$ (C1, C4) & 708 → 725 mm€ & 1 \\
\rowcolor{tceesalvia}\multicolumn{4}{l}{\emph{Trabajo, costes y precios}} \\
CLU nominales, 2025 & $\gamma_W-\gamma_{\mathrm{PAT}}=4{,}4-0{,}0$ (C3) & 4,4~\% & 2, 4 \\
Margen, 2025 & $\gamma_P-\gamma_{\mathrm{CLU}}=3{,}0-4{,}4$ (C3, C5) & $-1{,}4$ & 4 \\
\autor{Okun}, 2014-2019 y 2022-2025 & $\overline{\Delta U}/(\overline{\gamma_Y}-\overline{\gamma_{Y^p}})$: $-2{,}0/1{,}7$ y $-1{,}1/1{,}2$ (C1, C3) & $-1{,}2$ y $-0{,}95$ & 2 \\
Vacantes por parado & $0{,}70/14{,}1$ y $0{,}90/10{,}5$ (C3) & 0,05 → 0,09 & 2 \\
Diferencial de inflación, sep. 2026 & general $5{,}0-3{,}8$; subyacente $3{,}5-2{,}5$ (G6) & 1,2 y 1,0 & 4 \\
Último trimestre para la previsión & $(12\times3{,}7-32{,}5)/3$ (G6, C5) & $\approx4{,}0~\%$ & 4 \\
CLU acumulados, 2007-2013 & España frente a zona del euro (C3) & +1,2 / +12,1~\% & 3 \\
\rowcolor{tceesalvia}\multicolumn{4}{l}{\emph{Sector exterior}} \\
Ahorro menos inversión, 2025 & $25{,}0-21{,}7$ (C4) & 3,3 ($\approx$ CC 3,2) & 1, 3 \\
Sector privado, 2025 & capacidad de la nación $-$ saldo público $=4{,}3-(-2{,}4)$ (C7, C9) & +6,7~\% del PIB & 3 \\
Otras variaciones de la PIIN, 2025 & $\Delta\mathrm{PIIN}-\mathrm{CF}+\frac{\gamma^N}{1+\gamma^N}\mathrm{PIIN}_{t-1}=-2{,}7-5{,}0-2{,}2$ (C7, C9) & $-9{,}9$ puntos & 3 \\
Tipo de cambio dólar/euro & Brent en dólares entre Brent en euros (G7) & 1,05 → 1,17 & 3 \\
Exportaciones a EE.~UU. & suma de 12 meses, mar. 2025 → may. 2026 (G9) & $-10{,}8~\%$ ($\approx0{,}1~\%$ del PIB) & 3 \\
\rowcolor{tceesalvia}\multicolumn{4}{l}{\emph{Políticas}} \\
Tipo real, sep. 2026 & $2{,}50-5{,}0$ y $2{,}50-3{,}8$ (G6, G12) & $-2{,}5$ y $-1{,}3$ & 4 \\
Regla de \autor{Taylor}, España & $4+1{,}5\,(3{,}6-2)+0{,}5\times0{,}6$ (G6, C1) & 6,7~\% & 4 \\
ROE = ROA × apalancamiento & $14{,}2/0{,}96$ (C8) & 14,8 & 5 \\
Esfuerzo fiscal, 2025 & $\Delta b^{pe}=0{,}1-(-0{,}2)$ (C9, G13) & +0,3 & 6 \\
Saldo primario estabilizador, 2025 & $\frac{0{,}0249-0{,}0572}{1{,}0572}\times101{,}4$ (C9, G14) & $-3{,}1~\%$ del PIB & 6 \\
Ajuste déficit-deuda, 2025 & $\Delta d-(\text{déficit primario})-\text{bola}=-0{,}9-0{,}0+3{,}1$ (C9) & +2,2 & 6 \\
Bola de nieve, 2021-2025 & suma anual de $\frac{i-\gamma^N}{1+\gamma^N}d_{t-1}$ (C9) & $-31{,}4$ puntos & 6 \\
\bottomrule
\end{tabular}
\fuente{Cuadernillo del examen (C: cuadro; G: gráfico) y fórmulas de la guía (A.2 y A.3).}
\end{table}
